\documentclass[answers]{exam}
%\documentclass{exam}
\usepackage{../../mypackages}
\usepackage{../../macros}

\title{Fiche d'exercices - Chapitre 8 - Angles}
\author{N. Bancel}
\date{10 Décembre 2025}

% Mise en forme des solutions en bleu
\SolutionEmphasis{\color{blue}}
\renewcommand{\solutiontitle}{\noindent}

\begin{document}

\textbf{Collège Lycée Suger}
\hfill
\textbf{Mathématiques} \\

\textbf{Année 2025-2026}
\hfill
\textbf{Classe de 6ème} \\

{\let\newpage\relax\maketitle}


\begin{center}
\textbf{Essayez de vous forcer à bien rédiger en faisant ces exercices}
\end{center}

\section*{Calculs de mesures d'angles}

\textcolor{blue}{\textit{Dans votre cahier d'exercice, le titre de cet exercice doit être "Exercice 47 - Fiche d'exercices - Chapitre 8"}}

\begin{figure}[H]
    \centering
    \includegraphics[width=0.5\linewidth]{08_01.jpg}
    \captionsetup{labelformat=empty}
    \caption{\label{} Figure 2}
  \end{figure} 

  \begin{solution}
On cherche les mesures des angles $\widehat{MRN}$ et $\widehat{RMN}$.

\medskip
\textbf{1. Calcul de la mesure de l'angle $\widehat{MRN}$.}

\medskip

Je sais que les angles $\widehat{PRM}$ et $\widehat{MRN}$ sont adjacents et qu'ils forment ensemble l'angle $\widehat{PRN}$.  
Donc :
\[
\widehat{PRN} = \widehat{PRM} + \widehat{MRN}.
\]

Or les points $P$, $R$ et $N$ sont alignés, donc l'angle $\widehat{PRN}$ est un \textbf{angle plat}.  
Donc : $\widehat{PRN} = 180^\circ$. Et d'après la figure, $\widehat{PRM} = 65^\circ$.  
Donc :

\begin{align}
180^\circ &= 65^\circ + \widehat{MRN} \\
\widehat{MRN} &= 180^\circ - 65^\circ \\
\widehat{MRN} &= 115^\circ.
\end{align}

\medskip
\textbf{2. Calcul de la mesure de l'angle $\widehat{RMN}$.}

Je sais que les angles $\widehat{PMR}$ et $\widehat{RMN}$ sont adjacents et qu'ils forment ensemble l'angle $\widehat{PMN}$.  
Donc :
\[
\widehat{PMN} = \widehat{PMR} + \widehat{RMN}.
\]

Or le segment $[MP]$ est perpendiculaire au segment $[MN]$, donc l'angle $\widehat{PMN}$ est un \textbf{angle droit}.  
Donc : $\widehat{PMN} = 90^\circ$. Et d'après la figure, $\widehat{PMR} = 60^\circ$.  
Donc :

\begin{align}
90^\circ &= 60^\circ + \widehat{RMN} \\
\widehat{RMN} &= 90^\circ - 60^\circ \\
\widehat{RMN} &= 30^\circ.
\end{align}

\medskip

Ainsi,
\[
\boxed{\widehat{MRN} = 115^\circ \quad\text{et}\quad \widehat{RMN} = 30^\circ.}
\]
\end{solution}


\section*{Identifier une propriété}
\textcolor{blue}{\textit{Dans votre cahier d'exercice, le titre de cet exercice doit être "Exercice 33 - Fiche d'exercices - Chapitre 8"}}

\begin{figure}[H]
    \centering
    \includegraphics[width=0.5\linewidth]{08_02.jpg}
    \captionsetup{labelformat=empty}
    \caption{\label{} Figure 2}
  \end{figure} 

\textbf{Question bonus}

\setcounter{question}{2}

\begin{questions}
  \question Essayer d'écrire votre conclusion sous la forme d'un théorème mathématique rigoureux. Avec juste assez de mots. Ni trop, ni pas assez, pour dresser votre conclusion.
  \textit{En Mathématiques, cela ne s'appelle en réalité pas un théorème. C'est une "conjecture". C'est-à-dire quelque chose qui vous semble vrai, à partir d'observations. Mais qui n'a pas encore été démontré}
\end{questions}

\begin{solution}
  \begin{itemize}
  \item Si l’on trace un cercle ayant pour diamètre le segment $[AB]$, alors pour tout point $M$ placé sur ce cercle (différent de $A$ et de $B$), le triangle $AMB$ est rectangle en $M$.
  \item Quand on prend un point $M$ n’importe où sur un cercle de diamètre $[AB]$, l’angle $\widehat{AMB}$ est toujours un angle droit.
  \item Vocabulaire avec cercle circonscrit : Dans un triangle rectangle ABC, le centre du cercle circonscrit est le milieu de l'hypoténuse.
\end{itemize}
\end{solution}

\section*{Mesurer un angle en autonomie}

Tracer 3 points A, B, et C très ressérés les uns des autres. Mesurer les angles : $\widehat{ABC}$, $\widehat{ACB}$ et $\widehat{BAC}$.

\begin{figure}[H]
    \centering
    \includegraphics[width=0.3\linewidth]{08_06.jpg}
    \captionsetup{labelformat=empty}
  \end{figure} 

  \begin{solution}

    On ne va pas ici écrire la solution mais plutôt donner une indication. L'idée ici est de tracer toutes les droites et demi-droites nécessaires à la mesure de ces angles. Il faut que les traits dépassent les limites du rapporteur pour bien voir par transparence le trait de la droite en dessous de la graduation du rapporteur. Cet exercice est aussi un bon entraînement pour nommer les angles.

    \begin{figure}[H]
    \centering
    \includegraphics[width=0.7\linewidth]{08_07.jpg}
    \captionsetup{labelformat=empty}
  \end{figure} 

  Vous pouvez aussi vous entraîner à nommer les angles. Comment se nomme l'angle vert, comment se nomme l'angle bleu, puis l'angle rouge ?
  \end{solution}


\section*{Travailler la rédaction}

  \begin{figure}[H]
    \centering
    \includegraphics[width=0.6\linewidth]{08_05.jpeg}
    \captionsetup{labelformat=empty}
  \end{figure} 

  \begin{solution}
On cherche à déterminer la mesure de l'angle $\widehat{BFE}$ puis celle de l'angle $\widehat{EFD}$.

\medskip
\textbf{1. Calcul de la mesure de l'angle $\widehat{BFE}$.}

\medskip

On sait que les points $B$, $F$ et $D$ sont alignés.

\medskip

De même, les points $E$, $F$ et $C$ sont alignés.

\medskip

Les droites $(BD)$ et $(EC)$ se coupent en $F$, donc les angles $\widehat{BFE}$ et $\widehat{DFC}$ ont :

\medskip

\begin{compactitem}
  \item le même sommet $F$ ;
  \item leurs côtés portés par les mêmes droites $(BD)$ et $(EC)$, mais dans des sens opposés.
\end{compactitem}

\medskip

Les angles $\widehat{BFE}$ et $\widehat{DFC}$ sont donc \textbf{opposés par le sommet} et ont donc la même mesure. D'après la figure, $\widehat{DFC} = 44^\circ$, donc :
\[
\widehat{BFE} = 44^\circ.
\]

\medskip
\textbf{2. Relation entre les angles $\widehat{EFD}$ et $\widehat{DFC}$.}

\medskip

Les angles $\widehat{EFD}$ et $\widehat{DFC}$ :
\begin{compactitem}
  \item ont le même sommet $F$ ;
  \item ont un côté commun : la demi-droite $[FD)$ ;
\end{compactitem}

\medskip

Ils sont donc \textbf{adjacents}. De plus, les points E, F, C sont alignés donc l’angle $\widehat{EFC}$ est un \textbf{angle plat} :
\[
\widehat{EFC} = 180^\circ.
\]

Les angles $\widehat{EFD}$ et $\widehat{DFC}$ forment ensemble l’angle plat $\widehat{EFC}$, donc ils sont \textbf{supplémentaires}.

On peut donc écrire : 
\begin{align*}
\widehat{EFD} + \widehat{DFC} &= \widehat{EFC} \\
\widehat{EFD} + \widehat{DFC} &= 180^\circ \\
\end{align*}

\medskip
\textbf{3. Calcul de la mesure de l'angle $\widehat{EFD}$.}

\medskip

D’après la question 2, on a :

\[
\widehat{EFD} + \widehat{DFC} = \widehat{EFC}
\]
Or, d’après la figure, $\widehat{DFC} = 44^\circ$, et on sait aussi que $\widehat{EFC} = 180^\circ$ donc :
\begin{align*}
\widehat{EFD} + 44^\circ &= 180^\circ \\
\widehat{EFD} &= 180^\circ - 44^\circ \\
\widehat{EFD} &= 136^\circ.
\end{align*}

\medskip

Ainsi, on a trouvé :
\[
\boxed{\widehat{BFE} = 44^\circ \quad\text{et}\quad \widehat{EFD} = 136^\circ.}
\]
\end{solution}

\section*{Lire l'heure en classe}
\textcolor{blue}{\textit{Dans votre cahier d'exercice, le titre de cet exercice doit être "Lire l'heure en classe - Fiche d'exercices - Chapitre 8"}}

  \begin{figure}[H]
    \centering
    \begin{minipage}{0.45\linewidth}
        \centering
        \includegraphics[width=\linewidth]{08_03.jpg}
        \captionsetup{labelformat=empty}
        \caption*{Figure 1}
    \end{minipage}\hfill
    \begin{minipage}{0.45\linewidth}
        \centering
        \includegraphics[width=\linewidth]{08_04.jpg}
        \captionsetup{labelformat=empty}
        \caption*{Figure 2}
    \end{minipage}
\end{figure}

Une horloge est divisée en \textbf{12 grandes graduations}. 
Elle est consituée de deux aiguilles : 
\begin{compactitem}
    \item L'aiguille qui indique les heures (la plus courte)
    \item L'aiguille qui indique les minutes (la plus longue)
\end{compactitem} 

\vspace{1em}

Chaque grande graduation correspond à \textbf{1 heure}. \\
L'aiguille des minutes fait un tour complet en \textbf{60 minutes}. \\ 
L'aiguille des heures fait un tour complet en \textbf{12 heures}.  


\begin{questions}
    \question Combien de minutes y a-t-il dans une heure ?    
    \question Sachant qu'il y a 60 minutes réparties sur les 12 grandes graduations : à combien de minutes correspond 1 grande graduation d'heure ?
    \question On considère l'angle formé par une grande graduation et la suivante : quelle est la valeur de cet angle ? Justifier.

    \begin{figure}[H]
    \centering
    \includegraphics[width=0.3\linewidth]{08_03_grandangle.jpg}
    \captionsetup{labelformat=empty}
  \end{figure} 

    \question Une petite graduation correspond à 1 minute. On considère l'angle formé par une petite graduation et la suivante (donc un angle / intervalle d'une minute) : quelle est la valeur de cet angle ? Justifier.

        \begin{figure}[H]
    \centering
    \includegraphics[width=0.3\linewidth]{08_03_petitangle.jpg}
    \captionsetup{labelformat=empty}
  \end{figure} 

  \begin{questions}

\question Combien de minutes y a-t-il dans une heure ?

\begin{solution}
Une heure est composée de \textbf{60 minutes}.
\end{solution}

\question Sachant qu'il y a 60 minutes réparties sur les 12 grandes graduations : à combien de minutes correspond 1 grande graduation d'heure ?

\begin{solution}
On sait qu'il y a \textbf{60 minutes} réparties sur \textbf{12 grandes graduations}.

Donc une grande graduation correspond à :
\[
60 \div 12 = 5 \text{ minutes}
\]
\end{solution}

\question On considère l'angle formé par une grande graduation et la suivante : quelle est la valeur de cet angle ? Justifier.
  

\begin{solution}
On sait qu'un tour complet mesure $360^\circ$ et qu'il est divisé en \textbf{12 grandes graduations}.

Donc :
\[
360^\circ \div 12 = 30^\circ
\]

Ainsi, l'angle entre deux grandes graduations successives mesure \textbf{$30^\circ$}.
\end{solution}

\question Une petite graduation correspond à 1 minute. On considère l'angle formé par une petite graduation et la suivante : quelle est la valeur de cet angle ? Justifier.

\begin{solution}
On sait qu'un tour complet mesure $360^\circ$ et correspond à \textbf{60 minutes}.

Donc :
\[
360^\circ \div 60 = 6^\circ
\]

Ainsi, l'angle correspondant à \textbf{1 minute} mesure \textbf{$6^\circ$}.
\end{solution}

\question Dessiner \textbf{et calculer la mesure de l'angle formé par} les aiguilles pour indiquer les heures suivantes (on choisit toujours le plus petit angle). On considère que l'aiguille des heures \textbf{reste figée} sur le numéro de l'heure.

\textcolor{red}{\textbf{Essayez de rédiger vos calculs de manière claire et organisée. Le format idéal est une opération du type}
  \[
  \ldots \times \alpha + \ldots \times \beta = \ldots \quad \text{degrés}
  \]}

\begin{parts}

\part 4h

\begin{solution}
À 4h, l'aiguille des heures est sur le 4 et l'aiguille des minutes est sur le 12.

\medskip

L'écart est de \textbf{4 grandes graduations}.

\medskip

Or 1 grande graduation correspond à $30^\circ$, donc :
\[
4 \times 30^\circ + 0 \times 6^\circ = 120^\circ
\]

Ainsi, l'angle formé par les aiguilles est \textbf{$120^\circ$}.
\end{solution}

\part 4h30

\begin{solution}
À 4h30, l'aiguille des heures est sur le 4 et l'aiguille des minutes indique 30 minutes.

\medskip

On calcule l'angle de l'aiguille des heures :
\[
4 \times 30^\circ + 0 \times 6^\circ = 120^\circ
\]

On calcule l'angle de l'aiguille des minutes :
\[
0 \times 30^\circ + 30 \times 6^\circ = 180^\circ
\]

L'angle entre les aiguilles est donc :
\[
(0 \times 30^\circ + 30 \times 6^\circ) - (4 \times 30^\circ + 0 \times 6^\circ)
= 60^\circ
\]

Ainsi, l'angle formé par les aiguilles est \textbf{$60^\circ$}.
\end{solution}

\part 9h15

\begin{solution}
À 9h15, l'aiguille des heures est sur le 9 et l'aiguille des minutes indique 15 minutes.

\medskip

On calcule l'angle de l'aiguille des heures :
\[
9 \times 30^\circ + 0 \times 6^\circ = 270^\circ
\]

On calcule l'angle de l'aiguille des minutes :
\[
0 \times 30^\circ + 15 \times 6^\circ = 90^\circ
\]

La différence donne :
\[
(9 \times 30^\circ + 0 \times 6^\circ) - (0 \times 30^\circ + 15 \times 6^\circ)
= 180^\circ
\]

Ainsi, l'angle formé par les aiguilles est \textbf{$180^\circ$}.
\end{solution}

\part 8h35

\begin{solution}
À 8h35, l'aiguille des heures est sur le 8 et l'aiguille des minutes indique 35 minutes.

\medskip

On calcule l'angle de l'aiguille des heures :
\[
8 \times 30^\circ + 0 \times 6^\circ = 240^\circ
\]

On calcule l'angle de l'aiguille des minutes :
\[
0 \times 30^\circ + 35 \times 6^\circ = 210^\circ
\]

L'angle entre les aiguilles est donc :
\[
(8 \times 30^\circ + 0 \times 6^\circ) - (0 \times 30^\circ + 35 \times 6^\circ)
= 30^\circ
\]

Ainsi, l'angle formé par les aiguilles est \textbf{$30^\circ$}.
\end{solution}

\part 3h42

\begin{solution}
À 3h42, l'aiguille des heures est sur le 3 et l'aiguille des minutes indique 42 minutes.

\medskip

On calcule l'angle de l'aiguille des heures :
\[
3 \times 30^\circ + 0 \times 6^\circ = 90^\circ
\]

On calcule l'angle de l'aiguille des minutes :
\[
0 \times 30^\circ + 42 \times 6^\circ = 252^\circ
\]

L'angle entre les aiguilles est donc :
\[
(0 \times 30^\circ + 42 \times 6^\circ) - (3 \times 30^\circ + 0 \times 6^\circ)
= 162^\circ
\]

Ainsi, l'angle formé par les aiguilles est \textbf{$162^\circ$}.
\end{solution}

\part 5h58

\begin{solution}
À 5h58, l'aiguille des heures est sur le 5 et l'aiguille des minutes indique 58 minutes.

\medskip

On calcule l'angle de l'aiguille des heures :
\[
5 \times 30^\circ + 0 \times 6^\circ = 150^\circ
\]

On calcule l'angle de l'aiguille des minutes :
\[
0 \times 30^\circ + 58 \times 6^\circ = 348^\circ
\]

On calcule d'abord la différence :
\[
(0 \times 30^\circ + 58 \times 6^\circ) - (5 \times 30^\circ + 0 \times 6^\circ)
= 198^\circ
\]

On choisit le plus petit angle :
\[
360^\circ - 198^\circ = 162^\circ
\]

Ainsi, l'angle formé par les aiguilles est \textbf{$162^\circ$}.
\end{solution}

\end{parts}

\end{questions}


\end{questions}


\end{document}
